\section*{Chapitre \ref{ch:g2:measuring-time} --- Mesurer le temps}

\begin{solution}{exo:g2:measuring-time:1}
Le réveil, le petit-déjeuner, le jeu de l'après-midi, le dîner, le
coucher.
\end{solution}

\begin{solution}{exo:g2:measuring-time:2}
Instant~; \omterm{def:g2:measuring-time:duration}{durée}~; instant~; \omterm{def:g2:measuring-time:duration}{durée}.
\end{solution}

\begin{solution}{exo:g2:measuring-time:3}
La bougie bleue a duré plus longtemps. Les bougies ont commencé au
même instant, et celle qui brûle encore quand l'autre est finie est
celle qui dure le plus longtemps --- exactement la règle de la course
honnête.
\end{solution}

\begin{solution}{exo:g2:measuring-time:4}
À chaque fois qu'on le retourne, le sable met la \emph{même} \omterm{def:g2:measuring-time:duration}{durée} à
s'écouler par la taille étranglée. Un arbitre qui ne change jamais
est juste avec tout le monde.
\end{solution}

\begin{solution}{exo:g2:measuring-time:5}
Notre impression du temps est une mauvaise horloge~: les attentes
ennuyeuses se traînent et les moments amusants filent. La montre
compte régulièrement, quoi que Ben ressente~: l'attente a donc
vraiment été plus courte que le dessin animé --- elle a seulement
\emph{paru} plus longue.
\end{solution}

\begin{solution}{exo:g2:measuring-time:6}
Les réponses varient~; en général quelques secondes pour le nom et
une vingtaine de secondes pour les chaussures. Le comptage doit
rester calme et régulier.
\end{solution}

\begin{solution}{exo:g2:measuring-time:7}
Un éternuement~: des secondes. Une nuit de sommeil~: des \omterm{ex:g2:measuring-time:clock}{heures}.
Nouer un lacet~: des secondes. Une matinée d'école~: des \omterm{ex:g2:measuring-time:clock}{heures}.
\end{solution}

\begin{solution}{exo:g2:measuring-time:8}
À cinq \omterm{ex:g2:measuring-time:clock}{heures}, la petite aiguille pointe sur le $5$ (la grande droit
vers le haut)~; à cinq \omterm{ex:g2:measuring-time:clock}{heures} et demie, la petite aiguille pointe
\emph{à mi-chemin} entre le $5$ et le $6$ (la grande droit vers le
bas).
\end{solution}

\begin{solution}{exo:g2:measuring-time:9}
Fais courir chacun contre le même sablier~: à l'heure du bain, lance
le sablier et compte combien de retournements complets dure le bain~;
demain, fais de même pour la douche. Note les deux comptes (par
exemple «~bain~: $4$ retournements, douche~: $2$ retournements~»)~:
plus de retournements signifie plus long, et n'importe qui, avec le
même sablier, peut vérifier.
\end{solution}

\begin{solution}{exo:g2:measuring-time:10}
Tom est plus près de la vérité. Compter les secondes n'est honnête
qu'à un rythme calme et régulier --- environ un battement de cœur au
repos par nombre, ou «~un petit crocodile~» par compte. Enchaîner les
nombres à toute allure rend les secondes trop courtes.
\end{solution}
